%% --- 6.4 Solver Performance ---

\subsection{Solution Discovery Performance}
\label{sec:exp:solver}

To compare our discovery module against baseline systems, we report ADRS performance in Table~\ref{tab:solver}. Following~\citet{cemri2026adaevolve}, we report best-of-3 scores across seeds for each system.
For each seed, \sys{} selects the highest-scoring branch at the final search iteration.
For Sakana, ARC, AIR, DS, and \sys{}, each score is from independent canonical evaluator re-runs on the selected solver code, ensuring cross-system comparability.
Human, AdaEvolve, and EvoX scores (Gemini-3.0-Pro, best-of-3 runs) are from original publications~\citep{cemri2026adaevolve,liu2026evox}.

\begin{table}[t]
\centering
\small
\caption{Solution discovery performance on ADRS benchmark tasks (best-of-3 seeds). Sakana/ARC/AIR/DS/\sys{} scores are from independent canonical evaluator re-runs on submitted solver code. Human, AdaEvolve, and EvoX scores are from original publications. $^*$~indicates Gemini-3.0-Pro, all other systems use Gemini-3.1-Pro.}
\label{tab:solver}
\setlength{\tabcolsep}{4pt}
\resizebox{\textwidth}{!}{%
\begin{tabular}{@{}llcccccccccc@{}}
\toprule
\textbf{Task} & \textbf{Dir.} & \textbf{Human} & \textbf{AdaEvo$^*$} & \textbf{EvoX$^*$} & \textbf{Sakana} & \textbf{ARC} & \textbf{AIR} & \textbf{DS} & \textbf{\sys{}} & \textbf{\sys{}$^*$} \\
\midrule
Prism      & $\uparrow$   & 21.89  & \textbf{26.26}   & \textbf{26.26}   & \textbf{26.26}  & 26.25  & \textbf{26.26}  & \textbf{26.26}  & \textbf{26.26}  & \textbf{26.26} \\
Cloudcast  & $\downarrow$ & 626.24 & 637.10  & 623.69  & 627.11 & 690.37 & 734.28 & 620.09 & \textbf{618.08} & \textbf{618.08} \\
EPLB       & $\uparrow$   & 0.1265 & 0.1450  & 0.1453  & 0.1270 & 0.1266 & 0.1449 & 0.1284 & 0.1459 & \textbf{0.1461} \\
LLM-SQL    & $\uparrow$   & 0.6920 & \textbf{0.7520}  & 0.7300  & 0.7320 & 0.6757 & 0.7148 & 0.7307 & 0.7222 & 0.7115 \\
TXN        & $\uparrow$   & 2724.8 & 4310    & 4310    & 4184   & 3247   & \textbf{4311}   & 4286   & 3906   & 3861 \\
\bottomrule
\end{tabular}}
\end{table}

All systems match or exceed the human expert baseline on all five tasks, consistent with the observation of~\citet{cheng2025barbarians} that LLM-based agents rapidly converge to similar solution quality.
Sakana's BFTS produces competitive scores---matching the Prism ceiling and ranking second on LLM-SQL---even though its generated papers often misreport or cherry-pick these numbers (\S\ref{sec:exp:integrity}). The reported scores are from canonical re-evaluation of the deterministic best node, not the figures in Sakana's papers.
\sys{} exceeds the human baseline on every task and achieves the best overall score on Cloudcast and EPLB, demonstrating that verifiability does not sacrifice performance.

\subsection{Success Analysis} \label{sec:success-cases}

We highlight two top-scoring solutions whose code we inspected to verify algorithmic novelty.

\begin{figure*}[t]
    \centering
    % Left Subfigure: Cloudcast
    \begin{subfigure}[b]{0.48\textwidth}
        \centering
        \includegraphics[width=\linewidth]{figures/novel-idea-cloudcast-new.png}
        \caption{Cloudcast}
        \label{fig:novel_ideas_cloudcast}
    \end{subfigure}
    \hfill
    % Right Subfigure: EPLB
    \begin{subfigure}[b]{0.48\textwidth}
        \centering
        \includegraphics[width=\linewidth]{figures/novel-idea-eplb.png}
        \caption{EPLB}
        \label{fig:novel_ideas_eplb}
    \end{subfigure}
    \caption{Overview of the novel algorithmic pipelines generated by \sys{}. (a) For Cloudcast, the system integrates a continuous Fractional Multi-Commodity Flow LP relaxation with a robust Randomized Shortest Path Heuristic (SPH) ensemble, bridged through a deterministic log-transformed weighting mechanism. (b) For EPLB, the system employs a four-stage pipeline featuring composite-key topology snapping and zigzag GPU assignment.}
    \label{fig:novel_ideas}
\end{figure*}

For the Cloudcast task, a natural formulation is finding a minimum-weight directed Steiner tree to ensure that shared path prefixes minimize egress fees. \sys{} solves this by combining a Fractional Multi-Commodity Flow LP relaxation with an ensemble of Randomized Shortest Path Heuristics (SPH), as shown in Figure~\ref{fig:novel_ideas}(a). The LP relaxation produces fractional edge flows over the full network. To convert these into valid discrete paths, the solver applies a log-transformed weighting mechanism that biases the SPH ensemble toward high-flow edges, avoiding the disconnected subgraphs that pure randomized rounding produces. This approach achieves the best transfer cost among all systems (Table~\ref{tab:solver}), outperforming both the published human expert and leading agentic baselines.

On the EPLB task, algorithms are strictly evaluated on a combination of load-balancing efficiency and execution latency. To optimize both metrics, \sys{} adopts a topology-aware hierarchical placement strategy. As illustrated in Figure~\ref{fig:novel_ideas}(b), this pipeline progresses through four distinct stages: allocating experts to nodes, performing global replication, snapping to the topology, and finally assigning replicas to GPUs. While the global replication step intentionally relies on an iterative argmax update to preserve balancing quality, the system achieves microsecond-level execution through two major vectorized innovations. First, it utilizes a novel composite-key topology snapping mechanism which enables a single hardware-accelerated sort to replace slow Python-level comparators. Second, it distributes these sorted replicas using a fully vectorized zigzag assignment pattern computed in a single scatter operation. This hardware-aware approach achieves a highly competitive combined score (Table~\ref{tab:solver}) with 4.91ms execution latency.